\chapter{The Quality Inspector Portal}
\label{ch:qc}

\portalcover{Quality Inspector Portal}{For staff at the inspection benches}{secQC}

\section{Who this chapter is for}

You inspect finished garments and decide whether each one is good enough to send
to the customer. The system gives you three things the old paper method could
not:

\begin{bullets}
  \item The \textbf{right specification} for the garment in front of you, pulled
        up automatically from its tag.
  \item A picture of the design on which you can mark \textbf{exactly where} the
        fault is.
  \item A record of \textbf{every inspection ever done} on that garment, so you
        know if it has been through before.
\end{bullets}

\section{Your screen}

Open \menu{Quality Control}. There are three tabs across the top.

\begin{table}[H]
\centering
\small
\begin{tabular}{@{}p{32mm}p{106mm}@{}}
\toprule
\textbf{Tab} & \textbf{What it is for} \\
\midrule
Inspect          & Your working screen. Scan a garment and check it. \\
QC queue         & Everything currently sitting in quality control. \\
Defect analysis  & Where faults happen on each design, and which faults are most common. \\
\bottomrule
\end{tabular}
\end{table}

\shot{80-qc-inspect-empty}{The Inspect tab, waiting for you to scan a garment.}{fig:qc-empty}

\section{Inspecting one garment}

\begin{steps}
  \item Put the garment on the bench.
  \item Scan its tag. If your bench reader is set up, the number appears by
        itself; otherwise type it and press Enter.
  \item Press \btn{Look up}.
\end{steps}

The screen now shows everything about this garment.

\shot{82-qc-inspect-garment}{A garment ready for inspection. The design picture is
on the left; the order, customer, size and inspection history on the right.}{fig:qc-garment}

\begin{table}[H]
\centering
\small
\begin{tabular}{@{}p{34mm}p{104mm}@{}}
\toprule
\textbf{What you see} & \textbf{Why it matters} \\
\midrule
Design picture      & The correct design for this garment, front and back. \\
Colour and size     & Check the garment matches. \\
Order and customer  & Different customers accept different tolerances. \\
In QC since         & How long it has been waiting. \\
Previous failures   & If this is not the first time, look harder. \\
\bottomrule
\end{tabular}
\end{table}

\section{Marking a fault on the design}

This is the part that makes the retrofit operator's job possible.

\begin{steps}
  \item Look at the garment and find the fault.
  \item \textbf{Click the same place on the picture.} A small window opens.
  \item Choose the fault from the list, set how serious it is, and add a short
        note if it helps.
  \item Press \btn{Add defect}.
\end{steps}

\shotsmall{83-qc-defect-dialog}{Clicking the design opens this window. Choose what
is wrong and how serious it is.}{fig:qc-dialog}

A numbered marker now sits on the picture exactly where you clicked, and the
fault is listed beside it.

\shot{84-qc-defect-marked}{The fault is marked. Marker 1 on the design is the
same fault as item 1 in the list.}{fig:qc-marked}

\begin{table}[H]
\centering
\small
\begin{tabular}{@{}p{28mm}p{110mm}@{}}
\toprule
\textbf{Severity} & \textbf{Use it when} \\
\midrule
\texttt{CRITICAL} & The garment cannot be sold like this --- an open seam, a hole, a wrong label. \\
\texttt{MAJOR}    & A customer would notice and complain --- a broken stitch, a stain. \\
\texttt{MINOR}    & Cosmetic; a customer probably would not notice --- slight puckering. \\
\bottomrule
\end{tabular}
\caption{Choosing the right severity.}
\end{table}

\begin{tipbox}
\textbf{Use both views.} Press \btn{Front} or \btn{Back} above the picture to mark
faults on either side. If a fault has no particular place --- a measurement that
is out of tolerance, for instance --- use \btn{+ Add defect without a position}.
\end{tipbox}

\begin{tipbox}
Clicked in the wrong spot? Just click the marker itself and it disappears.
\end{tipbox}

\section{Passing or failing}

Two large buttons sit at the top right of the garment panel.

\begin{table}[H]
\centering
\small
\begin{tabular}{@{}p{38mm}p{100mm}@{}}
\toprule
\textbf{Button} & \textbf{What happens next} \\
\midrule
\btn{Pass} &
  The garment is cleared and can go to dispatch. \\
\btn{Fail --- send to retrofit} &
  A repair job is opened automatically, carrying every marker you placed.
  The garment then needs sending to the retrofit bench on a transfer note. \\
\bottomrule
\end{tabular}
\end{table}

\begin{stopbox}
\textbf{You cannot fail a garment without giving a reason.} If you press
\btn{Fail} with no faults marked, the system stops you. This is deliberate: the
retrofit operator has to know what to fix.
\end{stopbox}

After you decide, the screen offers \btn{Inspect the next garment} so you can
carry straight on.

\section{Passing a whole clean batch}

When you have checked a pile by hand and every piece is good, you do not have to
scan them one at a time.

\begin{steps}
  \item Scroll to \emph{Pass a whole batch} at the bottom of the Inspect tab.
  \item Bulk-read the whole pile into the scan box.
  \item Press \btn{Pass all scanned garments}.
\end{steps}

\begin{warnbox}
Every pass is still recorded individually against \emph{your} name. Only use the
batch button when you really have inspected every garment in the pile.
\end{warnbox}

Anything that could not be passed --- because it is not in quality control, for
example --- is listed afterwards so nothing is silently skipped.

\section{The QC queue}

The \menu{QC queue} tab shows everything in the department, oldest first.

\shot{81-qc-queue}{The QC queue. Click any row to inspect that garment.}{fig:qc-queue}

The four boxes at the top tell you at a glance:

\begin{bullets}
  \item how many are \textbf{awaiting inspection},
  \item how many have \textbf{passed} and are waiting to go to dispatch,
  \item how many have \textbf{failed} and are waiting to go to retrofit,
  \item how many \textbf{batches are inbound} from finishing.
\end{bullets}

Clicking any row takes you straight to the inspection screen for that garment.

\section{Defect analysis}

The \menu{Defect analysis} tab answers the question ``where do our faults keep
happening?''

\shot{85-qc-analysis}{Defect analysis. Every marker ever placed on this design is
shown on one picture, so patterns stand out.}{fig:qc-analysis}

\begin{bullets}
  \item The picture on the left shows \textbf{every fault ever recorded} on that
        design. A cluster of markers in one place usually means a machine or an
        operation needs attention, not a careless worker.
  \item The table on the right ranks faults by how often they occur.
  \item The bottom table shows inspector activity.
\end{bullets}

\begin{plainwords}
A high fail rate for one inspector does not automatically mean they are wrong ---
they may simply be checking the hardest styles. Use this table to start a
conversation, never to draw a conclusion on its own.
\end{plainwords}

\section{Receiving and sending batches}

Quality control receives from finishing and sends to either retrofit or dispatch.
Both work exactly as described in Section~\ref{sec:receive}
and Section~\ref{sec:dispatch}.

\begin{tipbox}
When sending passed garments to dispatch, the system checks each one really has
passed. Anything that has not is refused with a clear reason, so a failed garment
can never slip into a customer shipment.
\end{tipbox}
